\documentclass[a4paper,12pt]{article}
\usepackage{graphicx}
\usepackage[a4paper, top=0cm, bottom=2cm, left=2cm, right=2cm]{geometry} % Ajuste les marges
\usepackage{xcolor} % Pour ajouter des couleurs
\usepackage{hyperref} % Pour avoir des références colorées si nécessaire
\usepackage{eso-pic}         % Pour ajouter des éléments en arrière-plan

\usepackage[french]{babel}
\usepackage[T1]{fontenc}
\usepackage{mathrsfs}
\usepackage{amsmath}
\usepackage{amsfonts}
\usepackage{amssymb}
\usepackage{tkz-tab}

\usepackage{tikz}
\usetikzlibrary{arrows, shapes.geometric, fit}
\newcounter{correction} % Compteur pour les questions

% Définir la commande pour afficher une question numérotée
\newcommand{\question}{%
  \refstepcounter{correction}%
  \textbf{\textcolor{black}{Question \thecorrection :}} \ignorespaces
}
% Commande pour ajouter du texte en arrière-plan
\AddToShipoutPicture{
    \AtTextCenter{%
        \makebox[0pt]{\rotatebox{80}{\textcolor[gray]{0.6}{\fontsize{10cm}{10cm}\selectfont PGB}}}
    }
}
\begin{document}
\hrule % Barre horizontale
% En-tête
\begin{center}
    \begin{tabular}{@{} p{5cm} p{5cm} p{5cm} @{}} % 3 colonnes avec largeurs fixées
        Lycée Dindéfélo & \quad\quad\quad Correction Test 2 & ... Octobre 2026 \\
    \end{tabular}
    \\[-0.01cm] % Ajuster l'espace vertical entre le tableau et la barre
    \hrule % Barre horizontale
\end{center}
\begin{center}
    \textbf{\Large Correction - Calcul Littéral } \\[0.2cm]
    \textbf{\large Professeur : M. BA} \\[0.2cm]
    \textbf{Classe : 3 M1} \\[0.2cm]
\end{center}

\question (3 points)\\
\[
\begin{aligned}
A&=-2-6+2+6\\
&=-2+2-6+6\\
&=0+0\\
&=0\\
\end{aligned}
\quad\quad\quad
\begin{aligned}
B&=-12+11-8+12\\
&=-12+12+11-8\\
&=0+3\\
&=3\\
\end{aligned}
\quad\quad\quad
\begin{aligned}
C&=-9-15+9-21\\
&=-9+9-15-21\\
&=0-36\\
&=-36\\
\end{aligned}
\]
\[
\begin{aligned}
D&= 12 + 3 \times 5\\
&= 12 + 15\\
&= 27\\
&\\
&\\
\end{aligned}
\quad\quad\quad
\begin{aligned}
E&= 40 - (8 + 3 \times 2)\\
&= 40 - (8 + 6)\\
&= 40 - 14\\
&= 26\\
&\\
\end{aligned}
\quad\quad\quad
\begin{aligned}
F&= 40 - [8 + 3 \times (2 -  12 + 3) \times 5]\\
&= 40 - [8 + 3 \times (-7) \times 5]\\
&= 40 - [8 - 105]\\
&= 40 - (-97)\\
&= 137\\
\end{aligned}
\]
\\[1cm]
\question (1,5 points)\\
\[
\begin{aligned}
G&=(-2)\times(-7)+5\\
&=14+5\\
&=19\\
&\\
&\\
\end{aligned}
\quad\quad\quad
\begin{aligned}
H&=-15-(5)\times(-9)+9\\
&=-15-(-45)+9\\
&=-15+45+9\\
&=30+9\\
&=39\\
\end{aligned}
\quad\quad
\begin{aligned}
I&=2-(-11)\times(7)-15+2-(-10)\times(4)\\
&=2-(-77)-15+2-(-40)\\
&=2+77-15+2+40\\
&=79-15+2+40\\
&=106\\
\end{aligned}
\]
\\[1cm]
\question (1,5 points)\\
\[\begin{aligned}
J &= -\frac{7}{5} + \frac{9}{10} - \frac{12}{5} + \frac{4}{10} \\
   &= \left(-\frac{7}{5} - \frac{12}{5}\right) + \left(\frac{9}{10} + \frac{4}{10}\right) \\
   &= -\frac{19}{5} + \frac{13}{10} \\
   &= -\frac{38}{10} + \frac{13}{10} \\
   &= -\frac{25}{10} = -\frac{5}{2} \\
\end{aligned} 
\quad\quad\quad 
\begin{aligned}
K &= -\frac{7}{3} \times \frac{5}{2} + \frac{35}{4} \times \frac{4}{6} \\
   &= -\frac{35}{6} + \frac{35 \times 4}{4 \times 6} \\
   &= -\frac{35}{6} + \frac{35}{6} \\
   &= 0 \\
   &= \\
\end{aligned}
\quad\quad\quad 
\begin{aligned}
L &= \frac{7}{3} \times \frac{5}{2} - \frac{35}{4} \times \frac{4}{6} \\
   &= \frac{35}{6} - \frac{35 \times 4}{4 \times 6} \\
   &= \frac{35}{6} - \frac{35}{6} \\
   &= 0 \\
   &= \\
\end{aligned}
\]

\end{document}